% Adapted from the manuscript's Discussion and Conclusion; no new empirical claims.
\chapter{Conclusion}
\label{ch:conclusion}

This project presents PHGST-DM for joint bottleneck forecasting in multi-entity production systems. It combines a heterogeneous process--logistics graph, process-informed supervision, and recurrent-pattern-guided delay-aware modelling. The predictor reports the bottleneck station, onset, duration, observable process cause, and remaining order time, connecting resource-level constraints with order-level production progress.

In the evaluated simulation, the model achieves event F1-scores of 80.4\%, 77.8\%, and 76.8\% at onset limits of 5, 10, and 15~min, respectively. These scores exceed the strongest baseline at each limit by 8.2--22.1 percentage points. Detected-event duration MAE is lower by 0.86--1.03~min, equivalent to a 31.6\%--37.5\% reduction. Process-cause accuracy and macro recall also exceed those of the compared baselines. Fine-tuning only the remaining-time head reduces its MAE to 9.60, 9.27, and 8.91~min while retaining the event and cause outputs.

The component studies support complementary roles for the proposed designs. Broader resource observations improve advance-warning coverage; process-cause supervision improves cause identification with little change in logged station F1; and temporal-pattern clustering and prototype matching strengthen the prediction of more distant events. The sensitivity pilot provides preliminary support for label relevance. These findings remain specific to the evaluated simulation; real-factory transfer and the production benefits of forecast-guided interventions require further validation.

Future work will focus on two connected directions. The first is transfer and validation beyond the simulator. Real-factory observations or synchronized digital-twin streams should be combined with the current Isaac Sim data to test sensor noise, incomplete event logs, maintenance periods, operating-speed changes, and production-policy changes. Independently reviewed labels, sensitivity analysis for persistence and threshold choices, rare-event probability calibration, and uncertainty estimates will be needed to determine whether the resource relations and process-informed bottleneck definitions remain reliable. These comparisons can also prioritize the most informative resource signals.

The second direction is forecast-guided operational support. The current predictor reports bottleneck location, onset, duration, cause, and remaining order time, but its outputs are not yet used to change production schedules, reallocate workers or AGVs, or adjust material-transfer orders. Future evaluation should select feasible interventions under resource-capacity, execution-delay, and safety constraints, then measure their effects on bottleneck duration, order completion, and throughput. This will establish a closed loop from prediction to decision, execution, and production response.
